% ======================================================================
\section{Limits, threats to validity and future work}\label{sec:limits}\label{sec:disclosures}

This section separates three things: what the approach itself cannot do, what the evaluation cannot show, and what
could make the reported numbers wrong. Section~\ref{sec:hard} explains the first in detail.

\subsection{Limits of the approach}
\begin{itemize}
  \item \textbf{Detector floor.} In the ceiling analysis of an earlier version, about one needed sound in five could
        not be found with today's open detectors, mostly because it is quiet under speech or music; in the final system
        six needed sounds per set are heard by no detector at all. No later rule can recover them.
  \item \textbf{Presence is not production.} The on-screen check confuses an object being visible with the object
        making the sound; every alternative tried traded one error for the other.
  \item \textbf{No audio stage tried here separates a right off-screen name from a wrong one} well enough to add hits
        without two to five wrong pictures each (Section~\ref{sec:hard-confirm}).
  \item \textbf{Offline only.} A 15-s clip needs about five minutes of GPU time (one NVIDIA H200) before the first picture
        is drawn, and each picture try takes about 20\,s. The system cannot caption live television.
  \item \textbf{English speech only.} The speech transcript, used only to ask the vision model whether people are reacting to a sound,
        comes from a
        Whisper model~\cite{radford2023whisper} run for English; speech in other languages gives no useful context.
\end{itemize}
The scope of the task is set by design rather than by a limit: no model is trained, speech and music are never drawn,
and one picture is shown per sound, without words (Section~\ref{sec:problem}).

\subsection{Limits of the evaluation}
\begin{itemize}
  \item \textbf{No viewer study.} The cost model assumes that a wrong picture harms a viewer half as much as a missed
        sound. No deaf or hard-of-hearing viewer has confirmed this, or that a generated picture is understood faster than
        a caption in brackets. The metric is therefore a model of the viewer, not a measurement;
        Figure~\ref{fig:beta} shows that, on point estimates, the main conclusion holds for any weight of a wrong picture
        below that of a miss.
  \item \textbf{One annotator.} All labels are the author's; inter-annotator agreement is unknown.
  \item \textbf{Strict ``wrong''.} It counts redundant and late pictures, and some ``wrong'' pictures are real
        off-screen sounds missing from the labels (Section~\ref{sec:strict}).
  \item \textbf{Pictures checked by the author only.} The picture result rests on one blind glance test by the author; the
        automatic picture check was validated on known cases, not against a human rater.
\end{itemize}

\subsection{Side check on fresh AudioSet clips}\label{sec:fresh}
After the system was frozen, the final system and direct audio-to-image generation were run once, unchanged, on 340 of
the 422 fresh AudioSet-Strong clips (Appendix~\ref{app:datahistory}); the other 82 were left out to save GPU time,
before any score was read. No rule was set on these clips. AudioSet has no needed or on-screen marks, so the scoring
rule was fixed before the run: a needed sound is a labelled sound of the 33 consequential classes of the AudioSet
rounds (alarms, sirens, horns, dogs, glass, gunshots, explosions and similar), with runs of one family less than 2\,s
apart joined; a picture of any other labelled non-speech sound is neither a hit nor wrong; hits use the window of
Section~\ref{sec:matching}, and a second picture of the same sound counts as wrong. Table~\ref{tab:fresh} gives the
result (\fp{benchmark/gold/fresh_side_by_side.json}). AudioSet labels only some sounds of a clip and does not mark
what is on screen, so right pictures of unlabelled sounds count as wrong and visible sounds count as needed; the table
is a rough side-by-side, not a main result.

\begin{table}[htbp]
\centering\small
\caption{Fresh AudioSet-Strong clips (340 clips, 65 needed sounds), scored once after the system was frozen. Wrong:
other sound / nothing / repeat. Cost: (4 $\times$ misses + 2 $\times$ wrong) / clips, lower is better. Difference:
the final system's cost minus this system's, per clip, with a 95\,\% interval from a paired clip bootstrap
(100{,}000 draws) and a two-sided $p$-value.}
\label{tab:fresh}
\resizebox{\linewidth}{!}{%
\begin{tabular}{lrrrrll}
\toprule
System & Hits & Wrong & F1 & Cost & Difference (final $-$ this) [95\,\% CI] & $p$ \\
\midrule
Show nothing & 0 / 65 & 0 & 0.00 & 0.765 & $+0.041\ [-0.065, +0.141]$ & $0.45$ \\
Direct audio-to-image & 21 / 65 & 70 (43/19/8) & 0.27 & 0.929 & $-0.124\ [-0.224, -0.018]$ & $0.025$ \\
\textbf{Final system} & 12 / 65 & 31 (19/12/0) & 0.22 & 0.806 & & \\
\bottomrule
\end{tabular}}
\end{table}

\subsection{Threats to validity}
\begin{enumerate}[label=\textbf{T\arabic*.}, leftmargin=*]
  \item \textbf{Development versus test.} Every decision was made on the development set; the test set was only
        scored. The test cost fell by about a third as much as the development cost (Figure~\ref{fig:progression}), so
        the test numbers, not the development numbers, are the estimate for a new video.
  \item \textbf{Small sets.} With 71 and 88 clips, one or two pictures can decide a development comparison: one wrong
        picture moves the development cost by 0.028 and one hit by 0.056 per clip. The 95\,\% intervals of the main comparisons are 0.5 to 1.2 per clip wide.
  \item \textbf{Label changes.} The author re-checked the visibility of 14 disputed development sounds blind (3 changed;
        0 of 20 random controls changed); test labels were never changed.
  \item \textbf{One baseline not run.} A video-conditioned generator raised during supervision (for example
        MiniMax-H3) creates new video; the per-sound score cannot evaluate it, so it was not compared.
  \item \textbf{Scored versus live system.} The scoring harness and the live pipeline differ in four small ways that do
        not change when a picture starts; one of them (the direction of the kinship rule, Figure~\ref{fig:pipeline}) was not re-scored
        (Appendix~\ref{app:config}).
\end{enumerate}

\subsection{Future work}
Each item addresses one of the limits above.
\begin{enumerate}
  \item \textbf{Measure the viewer's costs}: a study with 5--10 DHH viewers (``helped / hurt / neither'' per clip) that
        estimates the weight of a wrong picture instead of assuming it.
  \item \textbf{A second annotator} for 30 clips (the packet exists), to measure label agreement.
  \item \textbf{The presence-versus-production question}: a much larger vision model, a comparison of the frame at the
        sound's onset with a frame just before it, or an external box around the named object.
  \item \textbf{A detector that hears sounds under speech and music}, tested on fresh AudioSet clips such as those of Table~\ref{tab:fresh}, on which no rule in this report was set.
\end{enumerate}
